% Generated draft: requires mathematical and source review.
\documentclass[11pt]{article}
\usepackage[margin=25mm]{geometry}
\usepackage{fontspec}
\setmainfont{Latin Modern Roman}
\usepackage{amsmath,amssymb,mathtools,unicode-math}
\setmathfont{Latin Modern Math}
\usepackage{enumitem,needspace,xurl,hyperref}
\setlength{\parindent}{0pt}
\setlength{\parskip}{5pt}
\setlength{\emergencystretch}{3em}
\newcommand{\N}{\mathbb{N}}
\newcommand{\Z}{\mathbb{Z}}
\newcommand{\Q}{\mathbb{Q}}
\newcommand{\R}{\mathbb{R}}
\newcommand{\C}{\mathbb{C}}
\newcommand{\F}{\mathbb{F}}
\newcommand{\A}{\mathbb{A}}
\newcommand{\PP}{\mathbb{P}}
\newcommand{\E}{\mathbb{E}}
\newcommand{\eps}{\varepsilon}
\newcommand{\dd}{\,\mathrm{d}}
\begin{document}
\section*{UnsolvedMath ID: 801000}
\textbf{Source problem code:} catalogue entry
\subsection{Definitions and mathematical statement}
Let \(X\) denote the intended topological ambient space, and let \(\mathcal K\subseteq\mathcal P(X)\) denote the intended class of Kakeya sets. Here \(\mathcal P(X)\) is the set of all subsets of \(X\), and \(\operatorname{Bor}(X)\) is the smallest \(\sigma\)-algebra containing the open subsets of \(X\). A set \(B\) is a minimal Kakeya set if
\[
B\in\mathcal K
\quad\text{and}\quad
\forall C\subsetneq B,\ C\notin\mathcal K.
\]
The supplied question is whether
\[
\forall A\in\operatorname{Bor}(X),\qquad
A\in\mathcal K\ \Longrightarrow\
\exists B\subseteq A\,
\left(B\in\mathcal K\ \land\
\forall C\subsetneq B,\ C\notin\mathcal K\right).
\]
Only \(A\) is required to be Borel; neither \(B\) nor its proper subsets are subject to a stated regularity requirement. Minimality is with respect to all proper subsets, not merely Borel proper subsets.

The supplied record does not define a Kakeya set or specify the ambient space, dimension, directions, or segment-length convention. Thus this is a precise logical formulation relative to the intended \(X\) and \(\mathcal K\), whose geometric meanings remain unresolved. It is not a question quantified over arbitrary choices of \(X\) and \(\mathcal K\).
\subsection{Short English statement}
Does every Borel Kakeya set contain a Kakeya subset from which no points can be removed while retaining the Kakeya property? The contained minimal set need not be Borel. The ambient space and the precise Kakeya convention require clarification from the source.
\subsection{Sources}
\begin{itemize}
\item Supplied local UnsolvedMath registry record: permanent ID \(801000\), exact source problem code \texttt{catalogue entry}, title \emph{Open Problems in Mathematical Logic --- Problem 12}; upstream revision \texttt{d371ee917ac9647fbb06cd6fc529deec6f3dda82}.
\item The supplied record attributes the problem to Barmpalias, Gao et al., with proposed year 2026, and identifies the source list as \emph{Open Problems in Mathematical Logic} (2026), Problem 12: \url{https://arxiv.org/abs/2608.26628}. The supplied provenance reports access on 2026-10-16 and extraction by \texttt{arxiv\_latex}; these are supplied provenance claims, not independently verified access or source verification.
\item The registry labels the problem \emph{open}. This is a supplied status label, not an independently verified mathematical or literature assessment.
\end{itemize}
\end{document}
